\documentclass[11pt]{article}
\usepackage[margin=25mm]{geometry}
\usepackage{fontspec}
\setmainfont{Latin Modern Roman}
\usepackage{amsmath,amssymb,mathtools}
\usepackage{unicode-math}
\setmathfont{Latin Modern Math}
\usepackage{xurl,hyperref}
\hypersetup{colorlinks=true,urlcolor=blue}
\setcounter{secnumdepth}{0}
\setlength{\parindent}{0pt}
\setlength{\parskip}{5pt}
\setlength{\emergencystretch}{3em}
\newcommand{\sref}[2]{\hyperlink{source:#1:#2}{[S#2]}}
\newcommand{\cataloguescope}[1]{\paragraph{Recorded target.}#1}
\newcommand{\R}{\mathbb R}
\newcommand{\C}{\mathbb C}
\newcommand{\N}{\mathbb N}
\newcommand{\Z}{\mathbb Z}
\newcommand{\Q}{\mathbb Q}
\newcommand{\F}{\mathbb F}
\newcommand{\eps}{\varepsilon}
\begin{document}
\section{LOCC Dense-Codeability of the Four-Qubit W State}\label{statement}

\subsection{Definitions and mathematical statement}

\paragraph{Shared notation.}$\mathbb N=\{1,2,\ldots\}$; the standard number fields have their usual meanings. All additional conventions are stated locally.


\paragraph{Definitions and conventions.} A qubit has Hilbert space $\mathbb C^2$, and the normalized four-qubit state is $|W_4\rangle=\tfrac12(|1000\rangle+|0100\rangle+|0010\rangle+|0001\rangle)$. In the report's two-receiver protocol, two sending qubits are locally encoded and transmitted to their assigned receivers. The two receivers then decode by LOCC. A state is LOCC dense-codeable when the achievable classical communication rate exceeds the rate available from the transmitted systems without the shared entangled state; with two sending qubits the benchmark is two bits per use.

\paragraph{Source scope and qualifications.} This is the report's two-receiver LOCC question, not global decoding, a single-receiver protocol, or a teleportation question. Receiver partition and the report's encoding class remain part of the target.

\paragraph{Catalogue mathematical target.}
\paragraph{Catalogue wording retained for traceability.}
Is the so-called four-qubit $W$ state in the LOCC dense-codeable class?

\paragraph{Saved literature context.} The supplied catalogue's assessment dated 2026-08-21 records: The cited dense-coding literature develops capacity bounds and examples but does not settle whether the standard four-qubit W state is LOCC dense codeable under the report's two-receiver protocol. Later W-state communication results address related protocols, not this exact question. The assessment is retained as dated provenance; it is not a new comprehensive literature or open-status certification.
\subsection{Short English statement}
Can the four-qubit W state provide a dense-coding advantage when its two receivers communicate only classically?
\subsection{Research attempt}
\paragraph{Model, date and investigated scope.}
OpenAI GPT-6 Astra (configured model \texttt{gpt-6-astra}), ultra reasoning; generated 11 October 2026. Exact Gram-matrix obstruction for the four diagonal-Pauli W-state codewords; no answer to asymptotic LOCC dense-codeability.

\paragraph{Sources and literature check.}
The complete supplied definition, including source qualifications and the dated literature assessment, was reviewed. The calculation below is an elementary restricted attempt motivated by [S2]; it does not claim novelty or a fresh literature-status determination.

\paragraph{Context and investigated scope.} The complete supplied definition, including source qualifications and the dated literature assessment, was reviewed. The calculation below is an elementary restricted attempt motivated by [S2]; it does not claim novelty or a fresh literature-status determination.

\paragraph{Testing the phase-encoding route.} Label the two sender qubits first and put $Z=\operatorname{diag}(1,-1)$. For $u,v\in\{0,1\}$ set $|w_{uv}\rangle=(Z^u\otimes Z^v\otimes I\otimes I)|W_4\rangle$. Directly using the four disjoint computational basis terms gives
\[
 \langle w_{uv}|w_{u'v'}\rangle
 =\tfrac14\bigl(2+(-1)^{u-u'}+(-1)^{v-v'}\bigr).
\]
The overlap is zero exactly when both bits change. Hence an orthogonal subset of these four codewords contains at most two words: each word has just one orthogonal partner. No decoder, including a global one, can distinguish more than two of them perfectly in one use. To see the necessity, perfect discrimination by effects $E_m$ implies $E_m|w_m\rangle=|w_m\rangle$ and $E_m|w_j\rangle=0$ for $j\ne m$, forcing orthogonality.

\paragraph{Why this fails to settle dense-codeability.} Orthogonality alone does not supply the receiver-restricted decoder. A proposed local $X$-basis measurement would require checking that its outcome distributions distinguish the two words; no such decoder is asserted here. The proved statement is solely the global obstruction to a larger perfect phase code. Computational-basis measurements are even worse: they erase all four encoding phases and give exactly the same output distribution for every $u,v$.

\paragraph{Remaining gap.} The report permits more general local unitaries, receiver adaptation, and asymptotic decoding with small error. Restricting to diagonal Pauli encoders and demanding zero error can lose the desired advantage. This failed four-word construction gives neither a universal two-bit upper bound nor a rate strictly above two bits for the prescribed two-receiver partition.

\paragraph{Verification and adversarial audit.}
Checked all overlap cases and the distinction between perfect one-use discrimination and asymptotic capacity; removed an unproved LOCC decoding inference. Model self-audit only; no independent human proof checking. Prompts A and B were applied to this canonical record and its qualified source scope. Checked all overlap cases and the distinction between perfect one-use discrimination and asymptotic capacity; removed an unproved LOCC decoding inference. Retained unresolved status and zero original-target progress; General unitary block encodings and LOCC decoders have not been bounded or constructed.

Coordinator recomputed all phase-code overlaps and confirmed the size-two orthogonality obstruction. No local decoder or asymptotic capacity bound follows. This additional check was model-assisted, not independent human verification.

\paragraph{Tools, process and contributions.}
Python 3 was used for file preparation and structural checks; no numerical experiment or theorem prover was used. Contribution prompt A was followed by prompt B. GPT-6 Astra with ultra reasoning drafted the research and performed its stated self-audit; any additional coordinator contribution is identified in the verification and audit. A GPT-6-based Codex coordinator assembled the manuscript and checked repository integration; its exact variant and reasoning configuration were not exposed. No human mathematical contribution, independent human proof check or formal proof certification is supplied. Complete derivations and the consulted source references are included here; private reasoning transcripts are not published.

\paragraph{Final status and first remaining gap.}
\textbf{UNRESOLVED}. General unitary block encodings and LOCC decoders have not been bounded or constructed.

\section{Completion Estimate}
\noindent\textbf{COMPLETION: 0\%}
The calculation is an auxiliary or restricted attempt and establishes no newly verified portion of the full source target. General unitary block encodings and LOCC decoders have not been bounded or constructed.

\subsection{Sources}
\begin{itemize}
\item[\hypertarget{source:statement:1}{S1}] ULAM.AI and the UnsolvedMath Contributors, UnsolvedMath, catalogue entry. Catalogue curation under CC BY 4.0; definitions and English exposition added locally.
\url{https://huggingface.co/datasets/ulamai/UnsolvedMath}.
\item[\hypertarget{source:statement:2}{S2}] Entanglement and Decoherence: Mathematics and Physics of Quantum Information and Computation (2005), Oberwolfach Reports, DOI 10.4171/owr/2005/04, pp. 185--255; official report source and mathematical context.
\url{https://doi.org/10.4171/owr/2005/04}.
\item[\hypertarget{source:statement:3}{S3}] Official EMS Press PDF of the same Oberwolfach report; PDF pages 19--21, printed pages 203--205.
\url{https://ems.press/content/serial-article-files/45978}.
\item[\hypertarget{source:statement:4}{S4}] Bruss et al., Dense coding with multipartite quantum states; reference in the saved dated catalogue literature review.
\url{https://arxiv.org/abs/quant-ph/0507146}.
\item[\hypertarget{source:statement:5}{S5}] Yeo--Chua, teleportation and dense coding with genuine multipartite entanglement; reference in the saved dated catalogue literature review.
\url{https://doi.org/10.1103/PhysRevLett.96.060502}.
\end{itemize}
\label{end:statement}
\end{document}
